\documentclass[12pt]{article}

\newcommand{\duedate}{-----}
\newcommand{\assignment}{Lecture 13 Summary: Transformer Variants \& State-Space Models}

% Change the following to your name and UNI.
\newcommand{\name}{Aksel Kretsinger-Walters, adk2164}
\newcommand{\email}{}

% NOTE: Defining collaborators is optional; to not list collaborators, comment out the
% line below. Maximum of two collaborators per problem set.
%\newcommand{\collaborators}{Vig Vigerton (\texttt{UNI}), Alice Bob (\texttt{UNI})}
% No collaborators on PS 0

\makeatletter
\def\input@path{{../}{../../}{../../../}} % add as many parents as you need
\makeatother
\input{pset_template_NNDL.tex} %% DO NOT CHANGE THIS LINE

\begin{document}
\psetheader %% DO NOT CHANGE THIS LINE

\section{Overview}
This lecture surveys current Transformer research, focusing on efficiency, memory, and interpretability. Topics include linear Transformers, mixture-of-experts (MoE) architectures, and state-space models (SSMs) such as Mamba.

\section{Linear Transformers}
Linear Transformers (Katharopoulos et al., 2020) use a kernel trick to linearize the attention computation:
\[
		\text{LinearAttention}(Q,K,V) = \phi(Q)(\phi(K)^\top V)
\]
where $\phi(\cdot)$ is a feature map applied elementwise.
By computing $(\phi(K)^\top V)$ first, complexity drops from $\mathcal{O}(n^2)$ to $\mathcal{O}(n)$.

\subsection{Common Kernels}
Typical $\phi$ choices include:
\begin{itemize}
		\item \textbf{ReLU kernel:} $\phi(x)=\text{ReLU}(x)$
		\item \textbf{ELU kernel:} $\phi(x)=\text{ELU}(x)+1$
		\item \textbf{Softmax approximation:} $\phi(x)=\exp(x)/\sqrt{d_k}$
\end{itemize}

\subsection{Tradeoffs}
Linear Transformers offer strong computational and memory efficiency, especially for long sequences.
However, scaled dot-product attention remains more expressive and empirically stronger on many benchmarks.

\section{Other Transformer Variants}
\begin{itemize}
		\item \textbf{Sparse attention:} tokens attend to a subset of others (e.g., Longformer).
		\item \textbf{Low-rank approximations:} attention matrix projected into lower-dimensional space (e.g., Linformer).
\end{itemize}

\section{Mixture of Experts (MoE)}
MoE introduces sparsity in the MLP blocks of Transformers—only a subset of “experts’’ is active for each token.
A gating network assigns mixing proportions $p_k$ over experts:
\[
		E_c = \sum_k p_{ck} \exp\!\left(-\frac{\|t_c - o_{ck}\|^2}{2\sigma^2}\right)
\]
Training increases weights of experts performing better than average, leading to specialization.
Modern examples: \textbf{Mixtral 8×7B} and \textbf{Switch Transformer} (top-1 routing).

\section{State-Space Models (SSMs)}
SSMs model temporal dynamics explicitly:
\[
		x_t = A x_{t-1} + B u_t + w_t, \quad y_t = C x_t + v_t
\]
where $x_t$ is a latent state and $y_t$ the observation.
They capture long-term structure efficiently and operate in continuous time.

\subsection{Mamba (Gu \& Dao, 2023)}
Mamba combines SSM efficiency with attention-like adaptability:
\[
		h(t) = A(x_t)h(t\!-\!1) + B(x_t)x_t, \quad y(t) = C(x_t)h(t)
\]
It computes transition parameters $(A,B,C)$ as functions of input data, maintaining relevant context for long sequences.

\section{Comparative Insights}
\begin{itemize}
		\item \textbf{Mamba advantages:} linear scaling, strong on long or streaming sequences, efficient for time-series.
		\item \textbf{Transformer advantages:} richer relational modeling, parallel training, multi-modal and transfer-learning strength.
\end{itemize}

\section{Conclusion}
Transformers remain dominant for general tasks, but linearized attention, MoE routing, and state-space formulations like Mamba are redefining efficiency and scalability for long-sequence modeling.

\end{document}